\documentclass[11pt]{article}
\usepackage[a4paper, margin=2.0cm]{geometry}
\usepackage[T1]{fontenc}
\usepackage[utf8]{inputenc}
\usepackage{booktabs, array, amsmath, xcolor, enumitem}
\usepackage[colorlinks=true, linkcolor=blue!50!black, urlcolor=blue!50!black]{hyperref}
\newcolumntype{L}[1]{>{\raggedright\arraybackslash}p{#1}}
\setlength{\parskip}{4pt}
\sloppy

% Every number below comes from experiments/openfhe/energy_validation_tables/
% (generated by experiments/openfhe/make_energy_validation_tables.py from
% aggregate.py --scenario=energy_validation).
\newcommand{\tab}[1]{../../experiments/openfhe/energy_validation_tables/#1}
\input{\tab{numbers.tex}}

\title{Energy and memory: limited validation at fixed frequency\\[2pt]
\large Account before any complete energy/memory campaign}
\author{Boungo Tr\'esor}
\date{\today}

\begin{document}
\maketitle

\noindent\textbf{Purpose.} As requested in your review of the Standard campaign, before any complete energy/memory campaign: (A) a test of a fixed CPU frequency (boost disabled) on this laptop, with the frequency, the power state and the temperature verified throughout; (B) a limited validation on a few BFV, BGV and CKKS operations of both libraries, with gross package energy as the primary measurement, idle-corrected energy as a complementary estimate with a per-session baseline and its sensitivity, and memory reported as the observed peak process RSS. These results are kept separate from the latency campaign (different CPU configuration) and are not combined with it; the latencies recorded inside the energy invocations are not reported. No comparison between the libraries is drawn here: the purpose is to validate the measurement. Raw data: \texttt{experiments/results/raw/energy\_validation/}; all tables are generated from the data.

\paragraph{Summary.}
\begin{itemize}[leftmargin=*, itemsep=2pt]
  \item \textbf{A fixed frequency is possible and stable.} With boost disabled, the measured core frequency held at \freqAMains~MHz for 10~minutes on mains and on battery, at \tempABattery--\tempAMains\,\textdegree C instead of \tempAOn\,\textdegree C with boost.
  \item \textbf{Energy must be measured on mains only.} At the same frequency and load, the package drew \powDropA\% less on battery (\powABattery~W against \powAMains~W): the difference does not come from the core frequency. Every energy measurement is therefore made on mains; an invocation during which mains power is lost, even for one second, is set aside and measured again.
  \item \textbf{Gross energy is repeatable.} In two sessions (\nInvPerSession\ valid invocations each, boost never on, never off mains), the ratio of the two invocations of a cell lies within $\pm$5\% for \nRatioFive\ of \nCellsSess\ cell-sessions and within $\pm$10\% for \nRatioTen; equivalence within $\pm$10\% (TOST) is established for \nEquiv\ of \nCellsSess. Between sessions, \nAcrossFive\ of \nAcross\ cells agree within $\pm$5\%. The exceptions are explained in Section~\ref{sec:exceptions}: short windows (OpenFHE key generation), a disturbed start of the first session, and the dispersion of single invocations.
  \item \textbf{Idle-corrected energy is much more sensitive.} The session baselines agree (\baseOne\ and \baseTwo~W), but the corrected energy is only \shareMin--\shareMax\% of the gross energy, so it moves by \sensMinTwo--\sensMaxTwo\% when the lowest or highest idle reading of the clean session is used instead of the median. This supports gross energy as the primary measurement.
  \item \textbf{Observed peak process RSS is stable}: \rssMin--\rssMax~MB; the three invocations of a cell differ by at most \rssRangeMax~MB, and the two sessions by at most \rssAcrossMax~MB.
\end{itemize}

\section{Part A: fixed frequency}

\textbf{Procedure} (\texttt{experiments/tools/fixed\_freq\_test.sh}). No FHE code: the core used by the harnesses (core~0) was kept busy by a plain loop, and every second the script recorded the mains state, battery level, boost state, the core's effective frequency (from the hardware counters), the requested frequency, the CPU temperature, both fan speeds and the package energy counter. Phases: boost on, on mains (3~min, the configuration of the latency campaign); boost off, on mains (10~min); boost off, on battery (10~min). CPU: AMD Ryzen~9 6900HX, driver \texttt{amd-pstate-epp}, governor \texttt{performance}, platform profile \texttt{balanced}. Boost and the governor are restored on exit.

\begin{center}\small\input{\tab{fixed_freq.tex}}\end{center}

\noindent\textbf{Observations.} With boost off the frequency did not vary (all samples within 7~MHz) except one sample at the second after unplugging (\lowABattery\ sample below 3200~MHz). With boost on it varied between about 3.9 and 4.5~GHz at \tempAOn\,\textdegree C, which is the variability the fixed frequency removes. The package power at the same frequency and load differs between mains and battery; the cause was not investigated (presumably a battery policy of the platform acting on parts of the package other than the core frequency; not verified), but the consequence is clear: \textbf{energy measured on battery is not comparable with energy measured on mains}.

\section{Part B: limited validation}

\textbf{Cells.} BFV, BGV and CKKS at $N=8192$, category~1 (a configuration where the three schemes support every operation in both libraries); key generation, encryption, multiplication, relinearization; SEAL~4.1.2 and OpenFHE~1.5.1: 24~cells. Key generation has the boundary of the latency campaign (initialization and key generation).

\textbf{Protocol} (\texttt{experiments/openfhe/run\_energy\_validation.sh}): the dedicated energy/memory procedure of the SEAL campaign, run with boost disabled for the whole session. Per cell:
\begin{itemize}[leftmargin=*, itemsep=1pt]
  \item \emph{Energy}, two invocations: package energy counter (RAPL) read around the timed window only; 30 repetitions after 2 warm-up; the operation repeated in each window to fill about 50~ms (inner loop from the Standard campaign latencies; key generation: one per repetition); just before the repetitions, the harness's own idle reading (6~s settle, 4~s window). Gross energy per operation = package energy of the window divided by the operations in it.
  \item \emph{Idle-corrected energy}: gross energy minus baseline $\times$ window, per operation, with the \emph{session} baseline = median of the session's 48 per-invocation idle readings. Two session idle windows (30~s settle, then 60~s with nothing running, at the start and at the end of the session) give an independent check.
  \item \emph{Memory}, three invocations, one operation per repetition (10 after 2 warm-up): observed peak process RSS (\texttt{VmHWM}, read by the process itself).
\end{itemize}
The 120 invocations of a session run in randomized order (seed recorded); every invocation is validated fail-closed (repetitions complete, every result checked correct). The power state (mains, battery, boost, platform profile, package counter) and the core frequency and temperature are recorded every second.

\textbf{Sessions.} e1 on \dateOne\ from \firstOne\ to \lastOne; e2 from \firstTwo\ to \lastTwo, started automatically after a 10-minute pause. They are independent sessions (separate processes, order, baseline), but on the same night and close in time.

\begin{center}\footnotesize\setlength{\tabcolsep}{4pt}\input{\tab{sessions.tex}}\end{center}

\noindent Boost was never on (\nBoostOn\ samples) and mains power was never lost (\nOffMains\ samples off mains; \nSetAside\ invocations set aside). The median frequency over each invocation was \freqInvMin--\freqInvMax~MHz. The temperature stayed low (median \tempOne\ and \tempTwo\,\textdegree C).

\subsection{Gross package energy (primary)}

\begin{center}\footnotesize\setlength{\tabcolsep}{4pt}\input{\tab{gross.tex}}\end{center}
\noindent{\footnotesize mJ/op: mean of the two invocations. inv.\,2/1: ratio of the two invocations in the session; $^*$: equivalence within $\pm$10\% not established (TOST, Welch). e2/e1: ratio between sessions. Window: mean timed window per repetition.}

\subsection{Exceptions}\label{sec:exceptions}

\begin{center}\footnotesize\setlength{\tabcolsep}{3pt}\input{\tab{not_equivalent.tex}}\end{center}

\noindent\textbf{Observations.} The \nNotEquiv\ exceptions fall into three groups.
\begin{itemize}[leftmargin=*, itemsep=1pt]
  \item \emph{Short windows.} OpenFHE key generation (in two sessions for BFV, one for CKKS) has a window of about 18~ms, one key generation per repetition. Its dispersion is the highest of the table, which makes a $\pm$10\% equivalence hard to establish with 30 repetitions.
  \item \emph{Disturbed start of session e1.} \nHighOne\ per-invocation idle readings of e1 exceed \highFactor\ times the session median, at \highTimesOne; none in e2. The e1 start idle window (\idleStartOne~W) is also above its end window (\idleEndOne~W). Most of these readings fall in the first minutes of e1, when the laptop was in use: the operator was starting the second session in another terminal and progress-monitoring commands were being run. Three of the exceptions have one invocation in that interval (OpenFHE BFV encryption, BGV multiplication, BGV relinearization); BGV multiplication is the largest deviation of the validation (\acrossMin\ between sessions). This is an explanation consistent with the timing, not a verified cause. All invocations are reported; none was removed.
  \item \emph{Dispersion of one invocation.} OpenFHE BGV encryption in e1 (ratio 0.95) and SEAL BFV key generation in e2 (ratio 1.05) have a ratio within $\pm$5\%, but one of their two invocations has a relative standard deviation of 16--17\%, so the test cannot establish equivalence within $\pm$10\%.
\end{itemize}

\subsection{Idle-corrected energy (complementary) and memory}

\begin{center}\footnotesize\setlength{\tabcolsep}{4pt}\input{\tab{net_memory.tex}}\end{center}
\noindent{\footnotesize Net: gross minus the session baseline $\times$ window. Baseline sensitivity: relative change of e2's net energy when its lowest or highest per-invocation idle reading replaces the session median. Peak RSS: mean of the three memory invocations.}

\noindent\textbf{Observations.} The two session baselines agree (\baseOne\ and \baseTwo~W). In the clean session e2, the two idle windows (\idleStartTwo\ and \idleEndTwo~W) agree with the median reading, and using either of them instead changes the net energy by at most \sensWinMaxTwo\%. Using the lowest or highest single reading changes it by \sensMinTwo--\sensMaxTwo\%. In e1, the highest reading (\idleMaxOne~W, from the disturbed interval) would change it by up to \sensMaxOne\%, and its start idle window by up to \sensStartMaxOne\%; the median is not affected by these readings. Between the sessions, \netAcrossTen\ of \nAcross\ cells agree within $\pm$10\% on net energy (range \netAcrossMin--\netAcrossMax). The idle-corrected value is therefore usable only as a complement, with the baseline and its sensitivity stated, as you proposed.

\section{Proposal for the complete energy/memory campaign (for your approval)}

\begin{enumerate}[leftmargin=*, itemsep=2pt]
  \item \textbf{CPU configuration}: boost disabled for every energy/memory session, governor \texttt{performance}; frequency, temperature and power state logged every second and reported per session. Results kept apart from the latency results, which were obtained with boost on.
  \item \textbf{Mains only}: an invocation during which mains power is lost is set aside and measured again after mains has been back for 2~minutes (implemented and tested).
  \item \textbf{No use of the laptop during a session}, including monitoring commands; progress is read only after the session ends.
  \item \textbf{Measurements}: gross package energy per operation as the primary result; idle-corrected energy as a complement, with the per-session baseline (median of the session's own idle readings), the two session idle windows, and the sensitivity to the lowest and highest readings; observed peak process RSS from three invocations.
  \item \textbf{Key generation}: more repetitions for the cells whose window is shorter than about 25~ms (OpenFHE key generation), for example 100 instead of 30, so that their uncertainty is comparable with the other cells.
  \item \textbf{Scope and repetition}: the configurations of the Standard campaign, SEAL and OpenFHE in the same sessions, each cell measured in two independent sessions, if possible on different days. From these sessions, about one minute per cell and library, i.e.\ roughly 3.5~hours per session for the Standard scope.
\end{enumerate}

\section{Limitations}

One laptop and one configuration ($N=8192$, category~1); two sessions on the same night, 10~minutes apart. RAPL measures the whole package (all cores and the uncore), not the process alone, which is why the gross energy includes the machine's background consumption and why the idle correction is only an estimate. On AMD the counter is exposed through the \texttt{intel-rapl} interface. The fixed frequency (\freqAMains~MHz) is not the frequency of the latency campaign (boost, about 4.2--4.5~GHz), so energy and latency are not combined. The cause of the lower package power on battery was not investigated further; it is avoided by measuring on mains only.

\end{document}
